% called by main.tex
%
\section*{Abreviaturas y acrónimos}
\addcontentsline{toc}{section}{Abreviaturas y acrónimos}
\label{sec::acronimos}

\vspace{10 mm}
%Institutions
\begin{description}[align=right,labelwidth=2cm] 
\item [OMS] Organización Mundial de la Salud
\item [ODS] Objetivos de Desarrollo Sostenible
\item [PERCLOS] \textit{Percentage of Eye Closure}
\item [LSTM] \textit{Long-Short Term Memory}
\item [SBC] \textit{Single-Board Computer}
\item [EAR] \textit{Eye Aspect Ratio}
\item [MAR] \textit{Mouth Aspect Ratio}
\item [POSIT] \textit{Pose form Orthography and Scaling with ITerations}
\item [NHTSA] \textit{National Highway Traffic Safety Administration}
\item [EEG] \textit{Electroencephalography}
\item [EOG] \textit{Electrooculography}
\item [ELM] \textit{Extreme learning Machine}
\item [SVM] \textit{Support Vector Machines}
\item [ACES] \textit{Average Eye Closure Speed}
\item [PERLVO] \textit{percentage of mout large vertical open over time}
\item [PERTITL] \textit{percentage of head-tilt over time}
\item [PNS] \textit{Percentage of Non-Steering}
\item [SDLP] \textit{Standard Deviation of Lane Position}
\item [BN] \textit{Bayesian Network}
\item [DNN] \textit{Deep Neural Network}
\item [CNN] \textit{Convolutional Neural Networks} 
\item [RCAN] \textit{Residual Channel Attention Network}
\item [LCRN] \textit{Long-Term Recurrent Convolutional Network}
\item [FPS] \textit{Frames Per Second}
\item [HRNN] \textit{Hierarchical Recurrent Neural Network}
\item [YOLO] \textit{You Only Look Once}
\item [MOL-TCE] \textit{Mixed-effect Ordered Logic Time Cumulative Effect}
\item [CLAHE] \textit{Contrast Limited Adaptative Histogram Equalization}
\item [AUC] \textit{Area Under the Curve}
\item [AWG] \textit{American Wire Gauge}
\item [RAM] \textit{Random Access Memory}
\item [MAE] \textit{Mean Absolute Error}
\item [HOG] Histograma de Gradientes Orientados



\end{description}

\newpage
\textbf{Respecto al uso de Inteligencia Artifical}
En este trabajo se utilizo Gemini 3.8 Flash para la realización de tareas en múltiples secciones de la metodología
Se genero código con este modelo para disminuir tiempos en los algoritmos de detección de rostro y PERCLOS, detección de bostezo, en la detección de distracción y la modularización del algoritmo final y traslado de lenguaje a C++. También se utilizo como colaborador en generación de ideas en la sección de cálculo de orientación de cabeza. Finalmente fue utilizado para apoyar en la generación de imágenes representativas en esta tesis. 